\chapter*{Resumen}
\addcontentsline{toc}{chapter}{Resumen}

Las metodologías probabilísticas avanzadas para el análisis del riesgo de inundación ---generación estocástica de eventos, análisis multivariante de extremos, corrección de sesgo climático, downscaling híbrido--- permiten representar de forma más completa la variabilidad y la incertidumbre asociadas al riesgo de inundación que los enfoques deterministas clásicos basados en escenarios únicos de diseño. Sin embargo, su adopción operativa en proyectos reales de ingeniería hidrológica sigue siendo limitada. La hipótesis central de esta tesis es que esa barrera no es fundamentalmente científica, sino operativa: la fragmentación entre fuentes de datos, métodos estadísticos y modelos físicos impide integrarlos en flujos de trabajo reproducibles y transferibles.

Esta tesis desarrolla \pyhydra como núcleo metodológico y librería Python modular para integrar y operacionalizar metodologías avanzadas de análisis probabilístico del riesgo de inundación bajo incertidumbre hidrológica y climática. Ese núcleo se despliega y documenta dentro de \hydra, la plataforma de integración que agrupa la interfaz web, los notebooks, el entorno Docker y los servicios auxiliares de ejecución. En conjunto, \pyhydra y \hydra organizan el flujo completo ---desde la adquisición automática de datos y la generación estocástica de escenarios hasta la ejecución masiva de modelos hidrológicos e hidráulicos y la obtención de mapas de período de retorno basados en impacto--- de forma reproducible y transferible. Las metodologías operacionalizadas incluyen la generación estocástica de precipitación mediante NEOPRENE, librería que contiene el proceso NSRP puntual y su extensión multisitio STNSRP, la generación general de series mediante CoSMoS, el análisis multivariante de extremos con cópulas y métodos bayesianos, la corrección de sesgo de proyecciones CMIP6, el downscaling híbrido basado en MaxDiss y k-NN, y la automatización de modelos externos HEC-HMS, SWAT, SFINCS y HEC-RAS.

El caso del río Besaya en Los Corrales de Buelna constituye el origen conceptual de la línea de investigación: en ese trabajo se evidenció, para el caso analizado, que el período de retorno del forzamiento meteorológico no coincide necesariamente con el del impacto hidráulico simulado, principio que fundamenta la filosofía metodológica de \hydra y justifica la necesidad de evaluar la frecuencia directamente sobre los impactos. La validación del marco se realiza mediante casos que cubren tres dimensiones complementarias: inundación estocástica basada en impactos (Besaya, Sant Llorenç des Cardassar y Calle 30), análisis de extremos bajo incertidumbre (DANA de Valencia de octubre de 2024) e integración de cambio climático y datos globales (IAHR 2022, Lago Tanganica, países andinos, SIMPCCe y Atlas de Riesgo Climático de Panamá). En conjunto, estos casos ejercitan subconjuntos distintos de \pyhydra y validan \hydra como marco de transferencia.

La contribución de la tesis es la demostración de que es posible operacionalizar de forma reproducible y transferible la cadena completa del análisis probabilístico del riesgo de inundación bajo incertidumbre hidrológica y climática. \pyhydra no sustituye las metodologías existentes, sino que las integra en un núcleo común que facilita su aplicación sistemática en problemas reales de ingeniería, mientras que \hydra demuestra su transferencia como producto técnico navegable, ejecutable y documentado.

\textbf{Palabras clave:} HYDRA; análisis probabilístico de inundación; incertidumbre hidrológica; incertidumbre climática; marco reproducible; modelización estocástica; período de retorno basado en impacto.


\chapter*{Abstract}
\addcontentsline{toc}{chapter}{Abstract}

Advanced probabilistic methodologies for flood risk assessment --- stochastic event generation, multivariate extreme-value analysis, climate bias correction, hybrid downscaling --- allow a more complete representation of the variability and uncertainty associated with flood risk than classical deterministic approaches based on single design scenarios. Yet their operational adoption in real hydraulic engineering projects remains limited. The central hypothesis of this thesis is that this barrier is not fundamentally scientific but operational: the fragmentation among data sources, statistical methods and physical models prevents their integration within reproducible and transferable workflows.

This thesis develops \pyhydra as the methodological core and modular Python library for integrating and operationalising advanced probabilistic flood risk analysis methodologies under hydrological and climatic uncertainty. This core is deployed and documented within \hydra, the integration platform that brings together the web interface, notebooks, Docker environment and auxiliary execution services. Together, \pyhydra and \hydra organise the complete analysis workflow --- from automated data acquisition and stochastic scenario generation to large-scale hydrological and hydraulic model execution and the production of impact-based return period maps --- in a reproducible and transferable manner. The operationalised methodologies include stochastic precipitation generation through NEOPRENE, the library that contains the point NSRP process and its multi-site STNSRP extension, general stochastic series generation through CoSMoS, multivariate extreme-value analysis with copulas and Bayesian inference, CMIP6 projection bias correction, hybrid downscaling based on MaxDiss and k-NN, and the automation of external models HEC-HMS, SWAT, SFINCS and HEC-RAS.

The Besaya river case study in Los Corrales de Buelna constitutes the conceptual origin of the research line: that work evidenced, for the case analysed, that the return period of meteorological forcing does not necessarily coincide with that of the simulated hydraulic impact, a principle that underpins the methodological philosophy of \hydra and justifies the need to evaluate frequency directly on impacts rather than on the forcing variable. The validation of the framework is carried out through cases covering three complementary dimensions: impact-based stochastic flood assessment (Besaya, Sant Llorenç des Cardassar and Calle 30), uncertainty-aware extreme-value analysis (the October 2024 Valencia DANA), and the integration of climate change and global data (IAHR 2022, Lake Tanganyika, the Andean countries, SIMPCCe and the Panama Climate Risk Atlas). Together, these cases exercise different subsets of \pyhydra and validate \hydra as a transfer framework.

The contribution of the thesis is the demonstration that the complete chain of probabilistic flood risk analysis under hydrological and climatic uncertainty can be operationalised in a reproducible and transferable manner. \pyhydra does not replace existing methodologies; it integrates them within a common core that facilitates their systematic application to real engineering problems, while \hydra demonstrates its transfer as a navigable, executable and documented technical product.

\textbf{Keywords:} HYDRA; probabilistic flood analysis; hydrological uncertainty; climatic uncertainty; reproducible framework; stochastic modelling; impact-based return period.
